\begin{proof}[Proof of \cref{obj:lem:published-model}]
\begin{enumerate}
\item For a directed graph \(G'\) on \(V\), write
\[
\mathcal N^{\mathrm{in}}_i(G')
:=\{j\in V:(j,i)\in G'\},
\qquad
\mathcal N^{\mathrm{out}}_j(G')
:=\{i\in V:(j,i)\in G'\}.
\]
For a published schedule \(\vartheta=(G^\vartheta,c^\vartheta)\) and a binary assignment \(z\in\{0,1\}^V\), its polynomial outcome is
\[
Y_i^{\mathrm{pub}}(\vartheta,z)
:=
\sum_{J\subseteq\mathcal N^{\mathrm{in}}_i(G^\vartheta)}
c^\vartheta_{i,J}\prod_{j\in J}z_j.
\]
The published conventions of
\citep[Sections 3.1--3.2, Assumptions 1--2, equations (3.2)--(3.3)]{CortezRodriguezEichhornYu2023SNIPE}
give, for an interaction-order-one schedule,
\[
Y_i^{\mathrm{pub}}(\vartheta,z)
=
c^\vartheta_{i,\varnothing}
+\sum_{j\in\mathcal N^{\mathrm{in}}_i(G^\vartheta)}
c^\vartheta_{i,\{j\}}z_j,
\qquad
Y_{\max}(\vartheta)
=
\max_{i\in V}
\left(
|c^\vartheta_{i,\varnothing}|
+\sum_{j\in\mathcal N^{\mathrm{in}}_i(G^\vartheta)}
|c^\vartheta_{i,\{j\}}|
\right).
\]
These conventions permit a zero singleton coefficient on a graph arrow.

Let \(\theta^+\) denote the augmentation of an additive schedule \(\theta=(G,a,t,b)\):
\[
\theta^+:=(\widetilde G,c^{\mathrm{pub}}),
\qquad
\widetilde G:=G\cup\{(i,i):i\in V\}.
\]
Its singleton coefficients satisfy
\[
c^{\mathrm{pub}}_{i,\{j\}}
=
\begin{cases}
t_i,&j=i,\\
b_{ij},&j\ne i\text{ and }(j,i)\in G,\\
0,&j\ne i\text{ and }(j,i)\notin G.
\end{cases}
\]
The empty-set coefficient is \(a_i\). If
\(J\not\subseteq\mathcal N^{\mathrm{in}}_i(\widetilde G)\), then \(J\ne\varnothing\).
When \(J=\{j\}\), membership outside this neighborhood forces both
\(j\ne i\) and \((j,i)\notin G\), so its coefficient is zero.
When \(J\) is not a singleton, all singleton contributions in the coefficient construction vanish. Likewise, \(|J|>1\) forces \(J\) to be neither empty nor a singleton. Thus
\[
c^{\mathrm{pub}}_{i,J}=0
\quad\text{if}\quad
J\not\subseteq\mathcal N^{\mathrm{in}}_i(\widetilde G)
\ \text{or}\ |J|>1.
\]
This proves the interaction-order-one restriction.

Since \(G\) is off-diagonal,
\[
\mathcal N^{\mathrm{in}}_i(\widetilde G)
=\{i\}\,\dot\cup\,\mathcal N^{\mathrm{in}}_i(G),
\qquad
\mathcal N^{\mathrm{out}}_j(\widetilde G)
=\{j\}\,\dot\cup\,\mathcal N^{\mathrm{out}}_j(G).
\]
Splitting the coefficient sum at the own coordinate therefore gives
\[
|c^{\mathrm{pub}}_{i,\varnothing}|
+\sum_{j\in\mathcal N^{\mathrm{in}}_i(\widetilde G)}
|c^{\mathrm{pub}}_{i,\{j\}}|
=
|a_i|+|t_i|
+\sum_{j\in\mathcal N^{\mathrm{in}}_i(G)}|b_{ij}|.
\]
For \(\theta\in\mathcal M_{n,d}\), the coefficient-mass requirement in
\cref{obj:def:model} bounds each right-hand side by one. Taking the maximum proves
\[
Y_{\max}(\theta^+)\le1.
\]
The same disjoint neighborhood decompositions and the two degree bounds in
\cref{obj:def:model} give
\[
|\mathcal N^{\mathrm{in}}_i(\widetilde G)|
=1+|\mathcal N^{\mathrm{in}}_i(G)|\le d+1,
\qquad
|\mathcal N^{\mathrm{out}}_j(\widetilde G)|
=1+|\mathcal N^{\mathrm{out}}_j(G)|\le d+1.
\]
Every diagonal arrow belongs to \(\widetilde G\) by construction. These identities also cover \(d=0\), with empty off-diagonal neighborhoods and zero neighborhood sums.
% lean: CausalSmith.Experimentation.ThinnedgraphAdditiveRiskFrontier.augment_publishedClass

\item Conversely, let \(\vartheta=(G^\vartheta,c^\vartheta)\) satisfy the stated published restrictions and contain every diagonal arrow. Define its stripped schedule by
\[
\vartheta^-:=(G^-,a^-,t^-,b^-),
\qquad
G^-:=\{(j,i)\in G^\vartheta:j\ne i\},
\]
\[
a_i^-:=c^\vartheta_{i,\varnothing},
\qquad
t_i^-:=c^\vartheta_{i,\{i\}},
\qquad
b_{ij}^-:=c^\vartheta_{i,\{j\}}
\quad(i,j\in V).
\]
The graph \(G^-\) is off-diagonal. A source in a published in-neighborhood is either the recipient itself or a distinct source retained in \(G^-\). The analogous partition holds for recipients in an out-neighborhood. Consequently,
\[
\mathcal N^{\mathrm{in}}_i(G^\vartheta)
=\{i\}\,\dot\cup\,\mathcal N^{\mathrm{in}}_i(G^-),
\qquad
\mathcal N^{\mathrm{out}}_j(G^\vartheta)
=\{j\}\,\dot\cup\,\mathcal N^{\mathrm{out}}_j(G^-).
\]
Hence
\[
1+|\mathcal N^{\mathrm{in}}_i(G^-)|\le d+1,
\qquad
1+|\mathcal N^{\mathrm{out}}_j(G^-)|\le d+1.
\]
Cancelling one gives both degree bounds by \(d\). The published envelope bounds each recipient's coefficient mass, and splitting off the diagonal term yields
\[
\begin{aligned}
|a_i^-|+|t_i^-|
+\sum_{j\in\mathcal N^{\mathrm{in}}_i(G^-)}|b_{ij}^-|
&=
|c^\vartheta_{i,\varnothing}|
+\sum_{j\in\mathcal N^{\mathrm{in}}_i(G^\vartheta)}
|c^\vartheta_{i,\{j\}}|\\
&\le Y_{\max}(\vartheta)\le1.
\end{aligned}
\]
These are precisely the membership requirements in \cref{obj:def:model}, so
\[
\vartheta^-\in\mathcal M_{n,d}.
\]
% lean: CausalSmith.Experimentation.ThinnedgraphAdditiveRiskFrontier.strip_scheduleClass

\item Re-augmenting \(\vartheta^-\) recovers the published graph:
\[
G^-\cup\{(i,i):i\in V\}=G^\vartheta.
\]
Indeed, for \(j=i\), both graphs contain the arrow by the diagonal hypothesis; for \(j\ne i\), both contain it exactly when \((j,i)\in G^\vartheta\).

Write \(c^{(\vartheta^-)^+}\) for the coefficients of the re-augmented schedule. For every recipient \(i\) and every
\(J\subseteq\mathcal N^{\mathrm{in}}_i(G^\vartheta)\), we claim
\[
c^{(\vartheta^-)^+}_{i,J}=c^\vartheta_{i,J}.
\]
If \(J=\varnothing\), both coefficients equal \(a_i^-\).
If \(|J|>1\), both vanish by their order-one restrictions.
In the remaining case, \(J\ne\varnothing\) and \(|J|\le1\), so \(J=\{j\}\) for some \(j\). If \(j=i\), the re-augmented coefficient is
\(t_i^-=c^\vartheta_{i,\{i\}}\). If \(j\ne i\), the neighborhood condition gives
\((j,i)\in G^\vartheta\), hence \((j,i)\in G^-\), and the re-augmented coefficient is
\(b_{ij}^-=c^\vartheta_{i,\{j\}}\).
This proves recovery of every coefficient on a subset of the recipient's in-neighborhood.
% lean: CausalSmith.Experimentation.ThinnedgraphAdditiveRiskFrontier.augment_strip_equivalent

\item For any additive schedule, stripping its augmentation recovers its graph because
\[
(j,i)\in G\cup\{(k,k):k\in V\},\quad j\ne i
\quad\Longleftrightarrow\quad
(j,i)\in G,
\]
using the off-diagonal property of \(G\). The empty and own-singleton coefficient identities, followed by the off-diagonal singleton identity on graph arrows, give
\[
a^{(\theta^+)^-}=a,\qquad
t^{(\theta^+)^-}=t,\qquad
b^{(\theta^+)^-}_{ij}=b_{ij}
\quad\text{for every }(j,i)\in G.
\]
Together with graph recovery, these are the asserted additive-schedule correspondence.
% lean: CausalSmith.Experimentation.ThinnedgraphAdditiveRiskFrontier.strip_augment_equivalent

\item For any additive schedule, the order-one outcome formula and the disjoint augmented in-neighborhood give, for every \(i\in V\) and \(z\in\{0,1\}^V\),
\[
\begin{aligned}
Y_i^{\mathrm{pub}}(\theta^+,z)
&=
c^{\mathrm{pub}}_{i,\varnothing}
+\sum_{j\in\mathcal N^{\mathrm{in}}_i(\widetilde G)}
c^{\mathrm{pub}}_{i,\{j\}}z_j\\
&=
a_i+t_i z_i
+\sum_{j\in\mathcal N^{\mathrm{in}}_i(G)}b_{ij}z_j\\
&=Y_i(z).
\end{aligned}
\]
The second equality separates the own coordinate; every remaining source is distinct from \(i\) and lies on an arrow of \(G\), so its singleton coefficient is \(b_{ij}\).
% lean: CausalSmith.Experimentation.ThinnedgraphAdditiveRiskFrontier.pubOutcome_augment

\item Define the published population contrast by
\[
\tau^{\mathrm{pub}}(\vartheta)
:=
\frac1n\sum_{i\in V}
\left(
Y_i^{\mathrm{pub}}(\vartheta,\mathbf1)
-Y_i^{\mathrm{pub}}(\vartheta,\mathbf0)
\right),
\qquad
\mathbf1_j:=1,\quad \mathbf0_j:=0
\quad(j\in V).
\]
Applying the outcome identity at these two assignments gives
\[
\tau^{\mathrm{pub}}(\theta^+)
=
\frac1n\sum_{i\in V}\bigl(Y_i(\mathbf1)-Y_i(\mathbf0)\bigr)
=\tau(\theta).
\]
% lean: CausalSmith.Experimentation.ThinnedgraphAdditiveRiskFrontier.pubTte_augment

\item Let
\[
\operatorname{diag}(V):=\{(i,i):i\in V\},
\qquad
W_{ij}:=W(j,i)\quad(j\ne i),
\]
where \(W_{ij}\) is the audit mark for the arrow from \(j\) to \(i\).
For each assignment and audit realization, the retained graph and complete record are
\[
H:=\{(j,i)\in G:W_{ij}=1\},
\qquad
O:=\bigl(H,Z,(Y_i(Z))_{i\in V}\bigr).
\]
Define the record transformation
\[
\mathcal A(O)
:=
\bigl(H\cup\operatorname{diag}(V),Z,(Y_i(Z))_{i\in V}\bigr).
\]
On every off-diagonal pair,
\[
\mathbf1\{(j,i)\in\widetilde G\}\,W_{ij}
=
\mathbf1\{(j,i)\in G\}\,W_{ij}.
\]
Thus the retained off-diagonal graphs agree. The assignment coordinate agrees identically, and the outcome coordinates agree by the outcome identity already proved. Therefore, writing \(\widetilde O_{\theta^+}\) for the published record,
\[
\widetilde O_{\theta^+}=\mathcal A(O).
\]
On the space of records whose graph contains every diagonal arrow, the inverse transformation is
\[
\mathcal A^{-1}
\bigl(\widetilde H,Z,(Y_i(Z))_{i\in V}\bigr)
=
\bigl(\widetilde H\setminus\operatorname{diag}(V),
Z,(Y_i(Z))_{i\in V}\bigr).
\]
Adding and removing these known coordinates are measurable operations and are inverse to one another.
% lean: CausalSmith.Experimentation.ThinnedgraphAdditiveRiskFrontier.pubRecordOf_augment

\item Fix the common assignment-and-audit law satisfying
\cref{obj:ass:assignment,obj:ass:audit,obj:ass:design-independence},
and let \(\mathfrak C\) be a published schedule class such that
\[
\theta\in\mathcal M_{n,d}\quad\Longrightarrow\quad
\theta^+\in\mathfrak C.
\]
For any published schedule \(\vartheta\), its record under this channel is
\[
\widetilde O_\vartheta
:=
\left(
\{(j,i)\in G^\vartheta:j\ne i,\ W_{ij}=1\}
\cup\operatorname{diag}(V),
\ Z,\
(Y_i^{\mathrm{pub}}(\vartheta,Z))_{i\in V}
\right).
\]
As in \cref{obj:def:risk}, use an independent randomization seed with
\[
\xi\sim\operatorname{Unif}[0,1],
\qquad
\xi\mathbin{\perp\!\!\!\perp}(Z,W).
\]
For a measurable published estimator \(T^{\mathrm{pub}}\), define its squared risk and the corresponding additive-record estimator by
\[
\mathcal L^{\mathrm{pub}}(T^{\mathrm{pub}};\vartheta,q)
:=
\mathbb E\!\left[
\bigl(T^{\mathrm{pub}}(\widetilde O_\vartheta,\xi)
-\tau^{\mathrm{pub}}(\vartheta)\bigr)^2
\right],
\qquad
T(O,\xi):=T^{\mathrm{pub}}(\mathcal A(O),\xi).
\]
The latter estimator is measurable because \(\mathcal A\) is measurable. For every additive schedule, the record and target identities give the pointwise equality
\[
\bigl(T(O,\xi)-\tau(\theta)\bigr)^2
=
\bigl(T^{\mathrm{pub}}(\widetilde O_{\theta^+},\xi)
-\tau^{\mathrm{pub}}(\theta^+)\bigr)^2.
\]
Taking expectations under the same assignment, audit, and seed law therefore gives
\[
\mathcal L(T;\theta,q)
=
\mathcal L^{\mathrm{pub}}(T^{\mathrm{pub}};\theta^+,q).
\]
These expectations are understood as nonnegative extended integrals, so the identity also covers infinite risks.

For each \(\theta\in\mathcal M_{n,d}\), class inclusion now yields
\[
\mathcal L(T;\theta,q)
\le
\sup_{\vartheta\in\mathfrak C}
\mathcal L^{\mathrm{pub}}(T^{\mathrm{pub}};\vartheta,q).
\]
Taking the supremum over \(\theta\), and then using the infimum in
\cref{obj:def:risk}, gives
\[
R_{n,d}(q)
\le
\sup_{\theta\in\mathcal M_{n,d}}\mathcal L(T;\theta,q)
\le
\sup_{\vartheta\in\mathfrak C}
\mathcal L^{\mathrm{pub}}(T^{\mathrm{pub}};\vartheta,q).
\]
This holds for every measurable published estimator. Taking their infimum proves
\[
R_{n,d}(q)
\le
\inf_{T^{\mathrm{pub}}}
\sup_{\vartheta\in\mathfrak C}
\mathcal L^{\mathrm{pub}}(T^{\mathrm{pub}};\vartheta,q),
\]
which is the claimed minimax risk transfer under the common observation channel.
% lean: CausalSmith.Experimentation.ThinnedgraphAdditiveRiskFrontier.minimaxRisk_le_pubMinimaxRisk
\end{enumerate}
\end{proof}
